\documentclass[14pt]{article}

% --- Encoding / language (XeLaTeX-safe) ---
\usepackage{fontspec}
\usepackage[russian,english]{babel}
\defaultfontfeatures{Ligatures=TeX}
\setmainfont{Times New Roman}

% --- Math / theorem ---
\usepackage{amsmath}
\usepackage{amssymb}
\usepackage{amsthm}
\usepackage{stmaryrd}
\usepackage{breqn}
\usepackage{euscript}
\usepackage{mathrsfs}
\usepackage{enumerate}
\usepackage{empheq}

% --- Graphics / layout ---
\usepackage{graphicx}
\usepackage{wrapfig}
\usepackage{subcaption}
\usepackage{float}
\usepackage{xcolor}
\usepackage{geometry}

\geometry{
    left=3cm,
    right=3cm,
    top=2cm,
    bottom=2cm,
    bindingoffset=0cm
}

% --- Extra utilities ---
\usepackage{pifont}
\usepackage{lipsum}
\usepackage{todonotes}
\usepackage{csquotes}

% --- Bibliography ---
\usepackage[backend=biber,style=alphabetic]{biblatex}
\addbibresource{references.bib}

% --- Hyperref (must be loaded last) ---
\usepackage{hyperref}

% --- Custom imports ---
\input{notations}

% --- Safe image inclusion (shows empty figure if file missing) ---
\newcommand{\safeincludegraphics}[2][]{%
    \IfFileExists{#2}{%
        \includegraphics[#1]{#2}%
    }{%
        \textcolor{gray}{\rule{0.8\linewidth}{0.6\linewidth}}%
    }%
}

% --- Theorems ---
\newtheorem{proposition}{Proposition}
\newtheorem{theorem}{Theorem}
\theoremstyle{remark}
\newtheorem{remark}{Remark}

% --- Colors ---
\definecolor{darkred}{RGB}{139,0,0}

% --- Date helper ---
\newcommand{\todayYYMMDD}{\the\year-\the\month-\the\day}


\begin{document}

\hfill
% \todayYYMMDD
\begin{center}
% \section*{Курс ``МНИС 2025/25''. \\Эссе}
\textit{Social Preferences and Conflict: A Conceptual Framework and Empirical Analysis.}\\
% \textit{Крайний срок сдачи - 17 декабря 2024 г., 23:59}\\
\textit{Author: ;}\\
\textit{Supervisor: ;}\\
\textit{Date: .}
\end{center}

% \todo[inline, color=lightgray]{
% TODO:\\
% ---tbd conventions---\\
% \ding{51} be cool\\

% \bigskip
% ---definitely---\\

% \ding{55} add github\\
% \bigskip
% ---maybe---\\
% \ding{51} all \texttt{labels} used from?\\
% }
\tableofcontents
\newpage 

\section{Conceptual Framework: Social Preferences and Conflict}

Social preferences are understood as social concerns individuals have for fairness, reciprocity, and the welfare of others that go beyond their concerns for their material self-interest. These preferences contradict the traditional economic assumption that agents act solely for their own benefit. Behavioral economics and psychology emphasize that people have tendencies like altruism, inequity aversion and reciprocity that have a strong impact on individuals' economic and social decisions (Fehr \& Schmidt, 1999; Bowles \& Polania-Reyes, 2012). In the first part of this paper, we will define social preferences, and its components, theoretical origins, psychology, evidence from the empirical studies. In the second part of our paper, we will explore how the concept of social preferences can be applied to one topic in development economics -- military conflict. We will analyze several articles to understand how military conflict affects social preferences in the developing economies and examine whether there are any contradictions and gaps in the available literature, so to propose a new aspect of studying this topic.

\subsection{1.1 Defining Social Preferences and Their Components}

Social preferences are a broad set of behavioral tendencies that differ from pure self-interest. One important component is reciprocity, which has both positive and negative forms. Positive and negative reciprocity refers to rewarding people who behave pro-socially and punishing people who behave antisocially or violate norms. (Fehr \& Gachter, 2000). Altruism is another important aspect and refers to individuals who experience positive utility from making others better off even at a cost to themselves (Andreoni, 1989). A closely related idea is inequity aversion, first proposed by Fehr and Schmidt (1999), which claims that individuals dislike inequity and will ``pay'' some utility in order to return outcomes closer to equality. Such inequity-averse individuals feel distress when a counterpart obtains much more or much less than they themselves, and they engage in behaviors that attempt to change this imbalance, to redistribute resources. Also, the notion of fairness plays an important role in social references. People tend to prefer equitable distributions and would even reject unfair allocations when they know they will receive a decreased payoff, as evidenced in ultimatum games (Henrich et al., 2001).

In addition to these components, there are also social image concerns---acting to display a positive image and conformity to social norms, which entails behavior in accordance with the consensus of the group even without external enforcement (Henrich et al., 2001). According to studies, trust and reciprocity underpin cooperation in society, especially in societies with high levels of market integration (Fehr \& Hoff, 2011). These dimensions encompass the economic and social interactions that are influenced by social preferences, going beyond traditional economic models of purely rational decision-making.

Cultural and historical factors are also connected with social preferences. Social norms regulate trust, fairness, and cooperation and vary across societies based on history and institutional arrangements (World Bank Group, 2015). Individuals in societies characterized by high reciprocity like those with stronger community ties, weaker social hierarchy, more extensive communal traditions of mutual aid, demonstrate greater levels of reciprocity than those in societies with more transactional social interaction (Henrich et al., 2001).

\subsection{1.2 Theoretical Foundations and Main Contributors}

The notion of social preferences has been developed with contributions from different disciplines including economics, psychology and evolutionary biology. According to classical economic theory, individuals are traditionally assumed to be rational agents that act to maximize their self-interest disregarding social motives (Becker, 1976). However, evidence from experimental economics showed that human behavior systematically deviates from this assumption (Fehr \& Fischbacher, 2003).

Rabin (1993) in his work developed models that incorporated reciprocity and demonstrated that some people consider the intentions of others when making decisions. Inequity aversion was introduced by Fehr and Schmidt (1999), who argued that individuals experience aversion from income differences. Bolton and Ockenfels (2000) proposed another model by focusing on relative payoffs rather than on absolute fairness.

Behavioral economists Ernst Fehr, Samuel Bowles, and George Loewenstein have elaborated on this idea, incorporating psychological insights. Research in evolutionary biology pointed out that social preferences could evolve to increase cooperation and group survival. Research using game theory (e.g. dictator and public goods games) shows that people take social preferences into account and do not always act out of self-interest (Henrich et al., 2001). The role of economic development was also examined as a factor that can affect preferences, with market integration found to promote fairness norms (Fehr \& Hoff, 2011).

\subsection{1.3 Psychological Foundations and Deviations from Traditional Economic Assumptions}

The psychological aspect of social preferences challenges traditional economic models by showing the role of emotions, cognitive biases, and social conditioning. An important foundation of social preferences is empathy, the ability to understand and experience the emotions of others. Neuroscientific research shows that witnessing another in pain stimulates similar brain areas as personal suffering, driving altruistic action (Batson, 1991). The other key driver is social norms and cultural transmission. Individuals learn social norms from childhood and from the environment in which they live and operate and tend to use social norms to guide their behavior even while they remain anonymous (Fehr \& Gächter, 2002).

Another key driver for social preferences is a strong reciprocity---the tendency to help if helped, to punish if harmed, regardless of whether it is costly to do so---which has been recognized to be a major element affecting social preferences. In contrast to the process of instrumental reciprocity where individuals act positive to uphold relationships, strong reciprocity is independent of long-term payoffs and conditional cooperation is often based on moral incentives (Fehr \& Fischbacher, 2003).

\subsection{1.4 Empirical Evidence Supporting and Challenging Social Preferences}

Many empirical studies seem to validate the existence of social preferences and its importance for decision-making. Henrich et al. (2001) conducted experiments on 15 small-scale societies, showing the striking variation occurring in fairness norms and levels of cooperation that were often linked to the extent of market integration. To test this, Fehr and Hoff (2011) showed that preferences for both fairness and punishment increased with the degree of market involvement, further supporting the notion that economic institutions shape social preferences.

However, the robustness of social preference models is challenged by some researchers. The observed behavior in experiments may be strategic rather than social (List, 2007). Others question if social preferences are indeed fixed traits or state-based behaviors shaped by situational factors. (Bowles \& Polanía-Reyes, 2012). Also, studies suggest that financial incentives can either reinforce or undermine social preferences, depending on how they will be structured (World Bank Group, 2015).

Despite these debates, there is a clear evidence that social preferences significantly influence economic behavior.

\subsection{1.5 Defining Conflict}

In Part 2 of this paper we will explore how the notion of social preferences could be related to and enforce one of the main topics in the development economics. We will explore the notion of conflict further, to see how armed confrontation in some low- and middle-income countries alters social preferences. Before exploring the selected literature sources, however, it is important to take a look as to how scholars and researchers define and distinguish conflict. This is critical for us so to further try to investigate the relationship of conflict and social preferences.

Conflict is one of the topics covered in developmental economics, and its definition varies a lot depending on the perspective taken by different scholars. Conflict is defined in the World Banks' World Development Report 2011 as the organized use or threat of violence by a variety of actors, including state military forces, nonstate armed groups, criminal gangs and other factions, to achieve political, economic or social objectives. This report highlights that conflict is not just the absence of peace but a cause and a consequence of underdevelopment. The report illustrates that conflict goes beyond traditional wars but also includes less violent forms of contention such as political local civil conflict that can take the form of social-collective violence due to political exclusion, economic disparity and social injustice (World Bank. 2011 p. 2).

In a different source of Garfinkel and Skaperdas ``Economics of Conflict'' the authors characterize conflict as a competition in which self-interested actors intentionally direct resources toward the accumulation of the means of coercion as opposed to purely productive ends. This definition emphasizes that conflict arises in environments where rights to property are imperfectly enforced, and people are permitted to produce goods but also to appropriate. Their framework differentiates winner-take-all clashes from more subtle exchanges involving negotiations or settlements under the constant threat of violence.

To add further depth to the discussion, the article ``War and Conflict in Economics: Theories, Applications, and Recent Trends'' by Kimbrough, Laughren, and Sheremeta describes conflict as a situation where agents intentionally take on personal costs that are socially inefficient in an attempt to obtain what they are calling ``wins and losses'' as private gains. By this definition, conflict is distinct from mutually beneficial exchanges in that the focus is on resources being directed away from productive uses to competitive, zero-sum activities (Kimbrough, E. p. 2).

Apart from definitions and typologies, these sources introduce the aspect of the fact that conflict drives social preferences, shaping human behavior. According to the World Development Report 2011, in environments of chronic violence, such social preferences as altruism, fairness, and trust in public institutions tend to break down. In these environments, participation in networks of crime---especially among youth---seems rational as a better alternative to the state-sponsored forms of activities. Insecurity and the resulting evolution of risk perceptions results into the establishment of cycles of violence and underdevelopment (World Bank. 2011 pp. 11--12). The report further notes that the disintegrating power of reciprocity and social cohesion by conflict not only leads to economic inefficiency but also strengthens the mistrust that encourages violence (World Bank. 2011 p. 18).

Garfinkel and Skaperdas also study the behavioral consequences of conflict in their work. They reach a different conclusion, however, arguing that the risk and destructiveness of open conflict leads risk-averse individuals to prefer negotiated settlements and cooperation to resorting to violence. These settlements help to mitigate the high costs of fighting by shifting social preferences toward safety and stability and trust. This adjustment in behavior plays a critical role in changing resource allocation and in the end in the long term economic developmental of the country. (Garfinkel, M. pp. 668--670).

The group of authors Kimbrough, Laughren, and Sheremeta in the third source we checked also show that exposure to conflict can alter social preferences, fostering a stronger inclination toward certainty and even increased reciprocity. Their findings suggest that the experience of conflict may lead individuals to adopt behavioral strategies aimed at avoiding the high risks and uncertainties of violent engagements and will be more focused on enforcing peaceful collaboration and fairness. (Kimbrough, E. p. 1).

We believe that the difference in all these approaches and conclusions proves the importance of reviewing the variety of sources as well as case studies so to get a better evidence of hoe social preferences can be changes in a situation of conflict and what are the potential implications for policymakers or practitioners in terms of potential for recovery and resilience in post-conflict societies.

\section{Literature Review}

Many countries suffered wars in the 20th and 21st centuries. This is the reason why the impacts of armed conflicts on individuals, on post-conflict social preferences are of big interest for policy makers.

At the beginning of part 2 of our paper we want to mention one fascinating article ``\textbf{War and Nationalism: How WW1 Battle Deaths Fueled Civilians' Support for the Nazi Party}'' (De Juan et al., 2024). This article describes powerful research examining how the exposure to wartime mortality can transform collective attitudes, a topic that is relevant to our paper as it concerns the impact of armed conflict on social preferences. While the article is really focusing more on nationalism rather than social preferences like reciprocity or inequity aversion it still demonstrates a key mechanism: the psychological process of witnessing death in the community is increasing ingroup cohesion, cooperation, and leading to outgroup derogation. The civilians who do not participate directly in the fight still experience the trauma of losing relatives and community members, which leads to the shift of their preferences, strengthening sense of group solidarity and fairness based on shared sacrifice. The authors show how in interwar Germany in the regions with high fatality rates people switched their electoral preferences towards right-wing nationalist parties, and this effect they can explain by war-induced changes to the way civilians understand their own collective identity and their responsibility towards others (De Juan et al., 2024, p. 148). The researchers digitized and machine-coded comprehensive ``loss lists'' covering 7.5 million German soldiers from World War I and developed a county-level casualty fatality rate. This measures the intensity of exposure to battle deaths in each region. They correlate this figure to the voting results of Reichstag elections that took place from 1920 to 1933. In doing so, they show how changes in war casualties are linked to changes in nationalist voting.

We believe the strength of this research is in the use of archival data, matched with detailed electoral records, that makes it possible to analyze the causal effect of war fatalities on political behavior. The study does a great job of tracking shifts in nationalist sentiment. However, since the link between war-related trauma and heightened nationalism could be misread and misused to rationalize nationalistic agendas, we believe there might be ethical concerns about the interpretation of results.

For the literature review in our paper, we have chosen sources that describe more recent armed conflicts in low- or middle-income countries. We want to find out how the concept of social preferences has been operationalized in these studies and to identify strengths, limitations, or ethical concerns in its application.

\subsection{2.1 Review of Articles on the Alteration of Social Preferences in Low- and Middle-Income Economies Experiencing Armed Conflict}

The article by Callen, Isaqzadeh, Long, and Sprenger \textbf{``Violence and risk preference: Experimental evidence from Afghanistan''} provides evidence of how exposure to violence in contexts of armed conflict leads to the change in economic decision-making, indirectly showing light on changes in social preferences. Though it is more focusing on risk preferences, the methodology and results are highly relevant to our paper's investigation of whether armed conflict can change social behavior and preferences such as reciprocity, altruism, inequity aversion, and fairness. This is the study that attempts to explore if and how exposure to trauma impacts risk-taking behavior based on the experiment in Afghanistan. The author introduces some of the findings that people exposed to fear, particularly in areas with a higher incidence of violent attacks have substantially higher risk tolerance under uncertainty, as part of their rationalization effort of how conflict can shape relevant behavioral outcomes The strengths of this approach are notable. The author also describes experiments that measure people's choice in a situation of certainty and in a situation of uncertainty. When individuals are asked to recall fearful events, they display what is known as a ``Certainty Premium,'' e.g. they value the risk-free outcomes more than those that delivered with risk (Callen et al., 2014, pp. 124--128). This suggests that psychological impact of violence is not just risk-aversion, it has a wider implication providing people with incentive not to act violently and to cooperate in social situations, shifting their preferences to reciprocity and fairness. There is an obvious strength in this approach. The actual conflict zone of Afghanistan increases the validity of the results obtained there, ensuring that the observed impacts aren't artificial laboratory results, but are real behavioral reactions to a dangerous setting. While there are strengths of this study, there are several limitations and ethical concerns raised as well. A major limitation is, that the research does not directly measure changes in social preferences we would want to focus at, like for example altruism or inequity aversion. As we seek to explore these dimensions more evidence is needed to assess direct causal relationships between violent conflict and social preferences. The ethics of conducting experiments within a conflict environment are fraught. The ethical side d this research is also questionable, as provoking individuals to recall traumatic events can also risk re-traumatization or psychological harm.

The article \textbf{``Can Conflict Affect Individuals' Preferences for Income Redistribution''} by Ori Shai (2023) discusses how armed conflict can affect social preferences, especially fairness and inequity aversion. The study reports that people in conflict-affected areas significantly increased their support for income redistribution in 2006, following the war between Israel and a Lebanese terror organization. This suggests that exposure to conflict can realign social preferences: individuals start favoring policies that would reduce income inequality (Shai, 2023, pp. 3072--3074). Such changes can be seen as reflecting increased inequity aversion and fairness, which could well justify this article as an important piece in our search for changing social preferences in areas impacted by violence. Shai works with data from the European Social Survey (ESS) and the International Social Survey Programme (ISSP). Comparing individuals' social preferences before and after the conflict, and comparing those in less affected regions (the control group) with those living in conflict-affected areas (the treated group), the study estimates the extent to which the conflict impacts social preferences (Shai, 2023, pp. 3075--3078). The findings are based on the analysis of responses to survey items such as: ``The government should take measures to reduce differences in income level,'' which shows the level of respondents' fairness concerns and income inequality aversion. It also attempts to explain the mechanisms that underlie the change by focusing on the way exposure to conflict reduces individuals' sense of personal agency. According to Shai, when people experience violence, they are less inclined to believe they can shape their lives, and this this why they support redistribution policies as a way of managing the threats of external shocks (Shai, 2023, p. 3073). One of the best thing about this study is its comprehensive approach: the author uses data from several dataset and leverages econometric methods so to maximize the validity of the results. Another good point is that this research is very much focused on redistribution of income, which is an important dimension of fairness, as well as inequity aversion, that it is related to the very topic of our paper.

The article \textbf{``Growing up in the Iran--Iraq war and preferences for strong defense''} by Farzanegan (2021) provides a valuable examination how exposure to armed conflict can reshape social preferences. Though study is focusing more on security and defense the examples it provides would be also applicable to the broader concepts of reciprocity, altruism and collective sacrifice. The authors assume that people who live through the war during the ``impressionable years'' (18--25 years old) have a lasting influence on their personalities, making them to prioritize a strong national defense over other social objectives. This finding reflects how traumatic experiences of being in a conflict can shape core values such as supporting collective welfare, which seems to be linked to altruism and cooperative behavior (Farzanegan, 2021, p. 1947). In putting this concept into practice, Farzanegan and Gholipour use data from the World Values Survey (WVS 2005--2009). The survey asks respondents to select the most important national goal from a list that would include economic expansion, environmental improvement, and strong defence forces. The researchers were able to confirm a significant preference for defense, which can be interpreted as a heightened social readiness to act in line with the collective good, even at personal cost, and is therefore indicative of an increase in altruistic and reciprocal behavior (Farzanegan, 2021, p. 1948). Looking further at social preferences, prioritizing preparedness for collective action, the research also explores respondents' attitude towards fighting for their country should another conflict arise. (Farzanegan, 2021, pp. 1949). Though the strategy used by this research seems to be very strong, the focus on defense-related outcome only is rather narrow, and unfortunately other components of social preferences are not measured directly.

In \textbf{``The struggle for Palestinian hearts and minds: Violence and public opinion in the Second Intifada''} (Jaeger et al., 2012), the authors take on the question of how exposure to violence in an armed conflict shapes the attitudes of the population. While it primarily looks at political preferences---e.g., support for negotiating with Israel and allegiance to specific political wings---it offers insights that can certainly be extrapolated to social preferences more generally. One noteworthy example in the article is the authors' finding that the rise in local Palestinian fatalities produces a temporary radicalization of public opinion. Research reveals a measurable decline support for negotiations during the immediate weeks following an incident of violence, indicative of a departure from moderate or collaborative attitudes (Jaeger et al., 2012, p. 360). While this research does not directly assess reciprocity, inequity aversion and altruism, the decline in willingness to engage in peace negotiations can be viewed as decline in cooperation or cooperative behavior which is a key aspect of social preferences. In the contexts where practices of fairness and mutual benefit are fundamental to social interaction, this shift away from negotiation may be symptomatic of other problems - breakdowns in trust and reciprocity, leading individuals to prioritize retribution and self-defense over negotiation. In this research the authors combine survey data from routine public opinion polls in the Palestinian territories with detailed fatality data from a human rights organization. This enables them to estimate the impact of armed conflict on political preferences by comparing numbers and timing of deaths against variations in measures of public sentiment. The study employs econometric techniques which allow to capture both the immediate and long-term effects of conflict. Time lags used in the analysis show that a radicalizing effect of violence is short-term, fading after about three months (Jaeger et al., 2012, p. 362). This temporal dynamic seems to be very important, as it implies that the impact of conflict on social preferences does not have to be permanent and that it opens the door to studying resilience and flexibility of social behavior under the conditions of stress. The diversity of estimation techniques employed by the authors strengthens the robustness of their findings and gives us extra confidence that violence has temporal and measurable impacts on public opinion (Jaeger et al., 2012, p. 363). We believe this study also has its limitation -- this s a very specific context of Palestinian--Israeli conflict in terms of long-term violence, its history, and socio/political complexities. While the insights from this study are extremely useful, we suspect they might not be fully applicable to other conflict settings where cultural norms, economic conditions, and institutional frameworks are very different.

The article \textbf{``Violent Conflict and Behavior: A Field Experiment in Burundi''} (Voors et al., 2012) provides an insightful empirical exploration of the effects of exposure to armed conflict on the social preferences of individuals in Burundi -- the country in East Africa. The study is primarily focusing on a things like altruism, risk-taking, and patience. One of the key findings show that that in the arears characterized by higher conflict-induced mortality individuals are more altruistic in their behavior toward neighbors. During the field experiments they have been asked to allocate different monetary payoffs between themselves and an anonymous partner. The social orientation score from purely selfish to fully altruistic shows that individuals living within communities affected armed conflict exhibited higher levels of prosocial behavior (Voors et al., 2012, p. 947). This interpretation suggests that high insecurity contexts may increase cooperative behavior to promote mutual support to the extent that even pure altruism can come into play. Conflict-induced shifts in social preferences are operationalized through series of carefully conducted field experiments. The authors created indices of conflict at the community level (such as the proportion of deaths due to attacks) and household-level victimization indices, and then they related these indices to the types of behavior in the course of controlled experimental games (Voors et al. (2012, p. 944). ) This experimental approach helped to capture changes in the level of concern for others and altruism in the contexts of high insecurity (Voors et al., 2012, p. 946). The clear benefit of this research is that in these field experiments described the data has been collected in the real post conflict environment. It allows the researchers to demonstrate a more credible relationship between violent conflict and changes in social preferences. The only limitation of this study is its focus on only one country ---Burundi, which limits the generalizability of the data to other conflict zones with different cultural, economic or institutional contexts.

In the paper \textbf{``The Age of Consequences: Unraveling Conflict's Impact on Social Preferences, Norm Enforcement, and Risk Taking''} (Coutts, 2024) another east African country is in the focus of research. The author explores how exposures to violent conflict shape social behavior in the long run examining the case of Rwandan genocide in 1994 during the civil war and characterized by extreme violence when large members of the Tutsi ethnic group have been murdered systematically by neighbors. The study focuses mainly on a series of experimental games that measure cooperation and reciprocity. It finds that the communities exposed to higher degrees of violence and hate-propaganda messages in the course of Rwandan genocide in 1994 show significantly higher cooperation in public goods games decades later (Coutts, 2024, pp. 49--51). This finding is very relevant for our papers as it not only proves that social preferences are shaped during an active period of violent conflict, but that the conflict still has its influence decades after it was over. Coutts operationalizes the idea of conflict-induced changes in social preferences through a set of field experiments. In Rwanda's Rusizi district, public goods games were conducted in 147 villages nearly 20 years after the genocide, with the monetary decisions of participants perfectly reflecting cooperative behavior (Coutts, 2024, pp. 51--53). To add to these fields experiments they also did a survey to measure trust, cooperation and perception of risk, which enriched the analyses and allowed to see the split based on the age at which the individuals experienced the conflict. Community members exposed to the genocide as children showed a significant increase in cooperation, while the older people showed a significantly stronger tendency towards norm enforcement, both punishment and reward (Coutts, 2024, 56). We believe this is greatly conducted research with the only limitation we can think about -- a very special historical and cultural environment of Rwanda, and extreme nature of genocide that might also question the external validity of the findings.

As well as in Coutts's study of post-conflict Rwanda the group of authors Firchow, P., Funk, J., \& Mac Ginty, R. in their work \textbf{``After war ends: Aid paradigms and post‐conflict preferences''} (Firchow et al., 2025) explore how the priorities of local population change long after the cessation of armed conflict. While the article focuses on aid preferences in post‐conflict Bosnia and Herzegovina, its insights are relevant to the wider issue of shifting social preferences in post‐conflict societies Thirty years after the war citizens of Mostar, Bosnia and Herzegovina, turned away from classical post‐conflict topics of interest --- transitional justice, security, and similar topics. They turned towards grounded demand for material development needs --- for better roads, jobs, and public space. Survey results show that decades after the war end signs of peace became more connected with the development outcomes rather than with abstract notions of fairness and justice (Firchow et al., 2025, p. 3). This shift towards more tangible, day‐to‐day benefits prove that people tend to value shared economic and social wellbeing ahead of punitive or retributive justice. In the current research on Bosnia and Herzegovina, post‐conflict social preferences have been measured through the Everyday Peace Indicators (EPI) methodology. In seven different neighborhoods in Mostar focus groups of locals were asked to: Identify the signs they are using to identify whether their community is at peace. Researched created a set of indicators to reflect the community's priorities. These indicators were also grouped into categories and assigned ``importance scores'' to measure their relative importance. In doing so, the researchers transformed high‐level social preferences into quantifiable, context‐relevant indicators that captures the dimensions of reciprocity for example ( through joint public projects enabling mutual support), altruism (evident in the generalized demand for better public services), and fairness (realized in the equitable distribution of aid and infrastructure) (Firchow et al., 2025, pp. 4--6). We believe and this scientific approach and the use of EPI methodology implemented by Coutts and other colleagues has a significant strength as it can translate abstract concepts into quantifiable indicators.

\subsection{2.2 Literature Contradictions and Gaps}

While the literature we cover in this review contains valuable information about armed conflict affecting social preferences, we can also observe some gaps and contradictions. Some studies like those by Callen et al. (2014) and Voors et al. (2012), claim that the experience of conflict promotes cooperative social behavior, which manifests in reciprocity and fairness in social interactions. Others, such as the work by Shai (2023) suggests that conflict may discourage inequity aversion and the call for redistributive policies. Some scholars like Jaeger (2012), mention a decline in support for negotiation in addition to move away from cooperative behavior in the regions of conflict. This is a very significant discrepancy in views that leads to the ongoing debate whether post trauma l people to seek security by being risk averse, reciprocal, cooperative and altruistic or whether it leads people to cope through risk-taking behavior and violence.

Also looking into the existing literature sources, we understand that modern research methods provide strong evidence that conflict influences risk-taking, cooperation, and even political attitudes. However, the idea of social preferences and its change in the armed conflict situation is not exactly in the focus, with insufficient number of studies looking at the components of social preferences- reciprocity, altruism, inequity aversion and fairness. Besides, the contextual specificity of many studies means that findings may be difficult to generalize across different conflict settings. These contradictions and unclarities highlight the demand for further studies that would take a more comprehensive and culturally sensitive approach to study how the armed conflict radically transforms social preferences.

\section{3.New Research Question}

As we can see many researchers explored various aspects of the relationship between armed conflict and human preferences including risk-taking behavior, political preferences, and some elements of social preferences. Still not too much is known about the dynamic evolution of preferences in a context of ongoing war as well as decades after the war conflict. However, it would be extremely important to understand which changes are taking place after the end of armed conflicts so we can develop effective remedies tore-construct societies. So, the new research topic we would like to explore would be \textbf{how long-term exposure to the armed conflict and post-conflict recovery processes change social preferences---specifically, reciprocity, inequity aversion, altruism and fairness. What would be the implications for the resilience of the country, post-conflict remedies and possible policy interventions in low- and middle-income countries.} Though Russia is a high -income country (according to the World Bank classification) the proposed question might have a huge significance nowadays for our country as well in view of Russia--Ukraine conflict when the policy makers in both countries would have to come up with post-conflict recovery measures. Insights into mechanism of military actions affecting social preferences can be key to designing remedies that promote community resilience.

To effectively study the proposed topic, we would use mixed methodology, combining quantitative and qualitative methods. By quantitative methods we understand, for example field experiments, organizing economic games, as was described in some studies we reviewed earlier so to measure dimensions of social preferences. Also, we would use longitudinal data collection- this method should be applied at multiple points during and after conflict, allowing us to track the evolution of social preferences. By qualitative component we understand the possibility to talk to focus groups so to collect personal experiences and perceptions of individuals affected by war.

The result of this future study would be very helpful for practitioners and policy makers -- it will provide them with implications on launching reconciliation programs for reinstating trust in post-conflict settings. If the findings of our future study would clearly show that the influence of conflict on social preferences is dynamic and long-lasting, policymakers could adopt some practics to track behavioral patterns over time. All these measures will help to promote post-conflict recoveries and sustainable development opportunities for low- and middle-income countries.

\section{Literature:}

\begin{enumerate}
    \item Andreoni, J. (1989). Giving with impure altruism: Applications to charity and Ricardian equivalence. \textit{Journal of Political Economy}, 97(6), 1447-1458. \url{https://doi.org/10.1086/261662}
    \item Batson, C. D. (1991). \textit{The altruism question: Toward a social-psychological answer}. Psychology Press. \url{https://doi.org/10.4324/9781315808048}
    \item Becker, G. (1976). \textit{The economic approach to human behavior}. University of Chicago Press. \url{https://casdev.unc.edu/ppeprogram/wp-content/uploads/sites/26/2020/02/Becker-The-Economic-Approach-to-Human-Behavior-1976.pdf}
    \item Bolton, G. E., \& Ockenfels, A. (2000). ERC: A theory of equity, reciprocity, and competition. \textit{American Economic Review}, 90(1), 166-193. \url{https://doi.org/10.1257/aer.90.1.166}
    \item Bowles, S., \& Polanía-Reyes, S. (2012). Economic incentives and social preferences: Substitutes or complements? \textit{Journal of Economic Literature}, 50(2), 368-425. \url{https://doi.org/10.1257/jel.50.2.368}
    \item Callen, M., Isaqzadeh, M., Long, J. D., \& Sprenger, C. (2014). Violence and risk preference: Experimental evidence from Afghanistan. \textit{American Economic Review}, 104(1), 123--148. \url{http://doi.org/10.1257/aer.104.L123}
    \item Coutts, A. (2024). The age of consequences: Unraveling conflict's impact on social preferences, norm enforcement, and risk-taking. \textit{Journal of Economic Behavior and Organization}, 218, 48--67. \url{https://doi.org/10.1016/j.jebo.2023.11.027}
    \item De Juan, A., Haass, F., Koos, C., Riaz, S., \& Tichelbaecker, T. (2024). War and nationalism: How WW1 battle deaths fueled civilians' support for the Nazi Party. \textit{American Political Science Review}, 118(1), 144--162. \url{https://doi.org/10.1017/S000305542300014X}
    \item Farzanegan, M. R., \& Gholipour, H. F. (2021). Growing up in the Iran--Iraq war and preferences for strong defense. \textit{Review of Development Economics}, 25(4), 1945--1968. \url{https://doi.org/10.1111/rode.12806}
    \item Fehr, E., \& Fischbacher, U. (2003). Why social preferences matter -- The impact of non-selfish motives on competition, cooperation, and incentives. \textit{The Economic Journal}, 112(478), C1-C33. \url{https://doi.org/10.1111/1468-0297.00027}
    \item Fehr, E., \& Gächter, S. (2000). Fairness and retaliation: The economics of reciprocity. \textit{Journal of Economic Perspectives}, 14(3), 159-181. \url{https://doi.org/10.1257/jep.14.3.159}
    \item Fehr, E., \& Hoff, K. (2011). Tastes, castes, and culture: The influence of society on preferences. \textit{World Bank Policy Research Working Paper No. 5760}. \url{https://ssrn.com/abstract=1910506}
    \item Fehr, E., \& Schmidt, K. M. (1999). A theory of fairness, competition, and cooperation. \textit{Quarterly Journal of Economics}, 114(3), 817-868. \url{https://doi.org/10.1162/003355399556151}
    \item Firchow, P., Funk, J., \& Mac Ginty, R. (2025). After war ends: Aid paradigms and post‐conflict preferences. \textit{World Development}, 189, 106916. \url{https://doi.org/10.1016/j.worlddev.2024.106916}
    \item Garfinkel, M. R., \& Skaperdas, S. (2007). Economics of conflict: An overview. In T. Sandler \& K. Hartley (Eds.), \textit{Handbook of Defense Economics} (Vol. 2, pp. 649-709). Elsevier. \url{https://doi.org/10.1016/S1574-0013(06)02022-9}
    \item Gintis, H., Bowles, S., Boyd, R., \& Fehr, E. (2005). \textit{Moral sentiments and material interests: The foundations of cooperation in economic life}. MIT Press.
    \item Henrich, J., Boyd, R., Bowles, S., Camerer, C., Fehr, E., Gintis, H., \& McElreath, R. (2001). In search of Homo economicus: Behavioral experiments in 15 small-scale societies. \textit{American Economic Review}, 91(2), 73-78. \url{https://doi.org/10.1257/aer.91.2.73}
    \item Jaeger, D. A., Klor, E. F., Miaari, S. H., \& Paserman, M. D. (2012). The struggle for Palestinian hearts and minds: Violence and public opinion in the Second Intifada. \textit{Journal of Public Economics}, 96, 354--368. \url{https://doi.org/10.1016/j.jpubeco.2011.12.001}
    \item Kimbrough, E. O., Laughren, K., \& Sheremeta, R. (2017, July 19). War and Conflict in Economics: Theories, Applications, and Recent Trends (MPRA Paper No. 80277). Retrieved from \url{https://mpra.ub.uni-muenchen.de/80277/}
    \item List, J. A. (2007). On the interpretation of giving in dictator games. \textit{Journal of Political Economy}, 115(3), 482-493. \url{https://doi.org/10.1086/519249}
    \item Rabin, M. (1993). Incorporating fairness into game theory and economics. \textit{American Economic Review}, 83(5), 1281-1302. \url{https://econweb.ucsd.edu/~jandreoni/Econ264/papers/Rabin\%20AER\%201993.pdf}
    \item Shai, O. (2023). Can conflict affect individuals' preferences for income redistribution? \textit{Journal of Population Economics}, 36, 3071--3096. \url{https://doi.org/10.1007/s00148-023-00963-z}
    \item Voors, M. J., Nillesen, E. E. M., Verwimp, P., Bülte, E. H., Lensink, R., \& van Soest, D. P. (2012). Violent conflict and behavior: A field experiment in Burundi. \textit{American Economic Review}, 102(2), 941--964. \url{https://doi.org/10.1257/aer.102.2.941}
    \item World Bank Group. (2015). Thinking socially, Chapter 2. In \textit{World development report 2015: Mind, society, and behavior} (pp. 51-70). World Bank. \url{https://doi.org/10.1596/978-1-4648-0342-0_ch2}
    \item World Bank. (2011). \textit{World Development Report 2011: Conflict, Security, and Development}. World Bank. \url{http://hdl.handle.net/10986/4389}
\end{enumerate}

\end{document}